% Continguts de la sessió sobre forces.
% Font: tex/draft/staging/Fuerzas_corregido.md (castellà)
% El \chapter el genera \sessio{} a partir de meta.tex; aquí només seccions.

\textit{
Com tantes altres vegades en física, ens toca començar aclarint un concepte del qual el llenguatge quotidià s'ha apropiat alegrement: la força.
}

\textit{
I, sentint-ho molt per alguns, per aquí no hi apareixeran ni Darth Vader ni Luke Skywalker. Tot i que la força que ens interessa també serveix per moure objectes.
}

\textit{
Comencem precisament per aquí.
}

\section{La pilota que no era on l'havíem deixat}

Imagina que estem mirant una pilota de futbol completament quieta. Apartem la mirada un moment i, quan tornem a mirar, la pilota és un metre més enllà.

No sabem què ha passat mentre no miràvem, però hi ha una cosa que podem afirmar amb bastant seguretat: en algun moment, la pilota ha hagut de canviar la seva velocitat.

Sembla una obvietat, però és important: no pots desplaçar-te d'un lloc a un altre sense haver-te mogut pel camí. Almenys, mentre ningú no demostri que les pilotes de futbol es teletransporten.

Així doncs, entre les nostres dues observacions ha passat alguna cosa. I això és precisament el que investigarem.

Perquè la pilota hagi acabat allà, el seu estat de moviment ha hagut de canviar en algun moment. Ha hagut de passar d'estar en repòs a moure's, encara que fos una estona. Va, diguem-ho sense embuts: en algun moment ha canviat la seva velocitat.

\textit{Anomenem \textbf{força} un tipus d'interacció que provoca un canvi en l'estat de moviment d'un objecte.}

Ens quedarem amb aquesta definició, tot i que una força també es pot fer servir per deformar les coses. Si esprems una pilota de goma entre les mans no la portes enlloc, però alguna cosa li estàs fent, i la pilota ho nota.

\section{Més a poc a poc, setciències}

Però espera.

Com diria l'Homer Simpson quan li intentaven explicar la tercera dimensió: «Més a poc a poc, setciències...».

Si poso mitja classe a estirar d'una corda en una direcció i l'altra mitja en la direcció contrària, és molt possible que la corda no es mogui ni un centímetre. Vol dir això que allà no hi ha forces?

I tant que n'hi ha. Que els ho preguntin als de la corda quan acabin, rebentats i amb les mans que els cremen.

Doncs què passa?

Que en mecànica clàssica el que decideix com canvia la velocitat d'un objecte no és cada força per separat, sinó la \textbf{força resultant}: la suma de totes les forces que actuen sobre ell.

Ja ho estàvem desitjant, escriure-ho, oi?

\[ \sum \vec F = m\,\vec a \]

I aquí convé dir les coses amb compte, perquè és on gairebé tothom rellisca:

\textit{Una força pot estar actuant sense que canviï el moviment de l'objecte, perquè altres forces poden estar compensant-la.}

Fixa't que \textbf{no} hem dit «sense que l'objecte es mogui». Un cotxe que circula per l'autopista a 120 km/h constants també té resultant nul·la, i ben cert que no està aturat. Resultant zero no vol dir quiet: vol dir que la velocitat no canvia.

\section{Les forces porten fletxa}

Hauràs vist que a la fórmula hi ha fletxetes a sobre d'algunes lletres. No són decoració.

Una força no queda descrita amb un número i ja està. Si t'empenyo, no és igual cap a on. Per descriure una força necessitem:

\begin{itemize}
    \item un \textbf{mòdul}: quant;
    \item una \textbf{direcció}: la recta al llarg de la qual actua;
    \item un \textbf{sentit}: cap a quin costat d'aquesta recta;
    \item i un \textbf{punt d'aplicació}: on s'aplica.
\end{itemize}
Per això les forces són \textbf{vectors}, i per això se sumen com a vectors i no pas com qui suma pomes. Dues forces de 100 unitats cadascuna poden donar 200, poden donar 0 i poden donar 141 i escaig, segons com les col·loquis. En el cas de la corda, mitja classe estava sumant fins a zero amb un entusiasme considerable.

La unitat de força al Sistema Internacional és el \textbf{newton}, amb símbol:

\[ \mathrm{N} \]

Endevina a qui volien homenatjar amb el nom.

I, com el joule, no és una unitat fonamental: es construeix a partir de les de sempre.

\[ 1\ \mathrm{N} = 1\ \mathrm{kg}\cdot\frac{\mathrm{m}}{\mathrm{s}^2} \]

Per fer-nos una idea: una poma d'uns 100 g pesa aproximadament un newton. Sí, una altra vegada la poma. La física té les seves manies, i sembla que les pomes tenen padrins.

\section{Quants tipus de força hi ha?}

Posem un altre exemple: una poma que penja de la branca d'un arbre, tranquil·lament, a la seva. De sobte es desprèn i cau a terra.

Aquest cas té molt a veure amb l'anterior. La poma estava quieta i ara es mou. Alguna cosa ha passat entremig.

I aquí convé fer una distinció petita però important: que la poma es desprengui de la branca \textbf{permet} que caigui, però la força que l'accelera cap a terra és la gravetat. La Terra feia tota aquesta estona que estirava d'ella, també mentre penjava de l'arbre. El que ha desaparegut no és la quietud: és la branca.

Així que aquesta força de la gravetat és precisament la primera que tenim per començar a investigar. Vegem què hi tenen a dir els científics.

Sabem que existeixen moltíssims tipus de forces a l'univers però, fins on sabem, totes es poden explicar a partir de només quatre \textbf{forces fonamentals}.

\section{«Fonamental», o com dir «fins aquí arribem»}

Els físics tenim una manera bastant curiosa de reconèixer que hem arribat al límit del que entenem.

En lloc de dir «fins aquí sabem explicar les coses, però no tenim ni idea de què hi ha més avall», fem servir una paraula molt més elegant: \textbf{fonamental}.

Així, una força fonamental és aquella que, fins on sabem, no podem explicar mitjançant altres forces més bàsiques. Queda bastant millor que reconèixer que, de moment, no hem aconseguit continuar estirant del fil.

Convé no perdre de vista aquest \textit{fins on sabem}. Fa cent anys la llista era una altra. D'aquí a cent, ja ho veurem.

Cadascuna d'aquestes quatre forces té les seves pròpies característiques i actua de manera diferent. Algunes només tenen efectes apreciables a distàncies diminutes, com les que trobem dins dels àtoms. Altres actuen a distàncies enormes, fins i tot entre galàxies.

I tenen una cosa bastant curiosa en comú: poden actuar sense que els objectes es toquin. Encara que ens sembli que empènyer una taula exigeix estar en contacte amb ella, a escala microscòpica el que passa és una interacció entre les partícules que formen els dos objectes. Però d'això en parlarem una mica més endavant.

\section{La gravetat, o atreure's sense tocar-se}

Bé, on ho havíem deixat? Ah, sí. La gravetat.

Estem acostumats a veure-la actuant a escala còsmica: manté els planetes a les seves òrbites, governa el moviment de les estrelles i contribueix a mantenir unides les galàxies. Però també està actuant contínuament al nostre voltant, amb bastant menys èpica.

A la física de Newton, la gravetat no és res més que l'atracció que experimenten dos objectes pel simple fet de tenir massa.

No importa si parlem de dues bales, d'una poma i la Terra o de dos planetes separats per milions de quilòmetres. Si tenen massa, s'atrauen gravitatòriament.

I ho fan sense necessitat de tocar-se.

També és, de bon tros, la més fluixa de les quatre. Ho dissimula bé perquè té un planeta sencer estirant cap al mateix costat, però qualsevol imant de nevera de dos euros és capaç de guanyar a la Terra i quedar-se amb el clip.

\textit{[Això sí: estem utilitzant el model de Newton. Més endavant, Einstein complicarà una mica les coses, com sol passar just quan ens pensàvem que ja ho havíem entès.]}

\section{L'electromagnetisme, o per què no travesses la cadira}

Bé, i quines són les altres tres forces fonamentals que coneixem?

A una escala molt més propera a la nostra vida quotidiana trobem la \textbf{força electromagnètica}. Segur que coneixes una de les seves propietats: les càrregues elèctriques de signe oposat s'atrauen i les del mateix signe es repel·leixen.

Aquesta força és responsable de l'estructura dels àtoms i de les molècules. Els electrons, amb càrrega negativa, són atrets pels nuclis atòmics, amb càrrega positiva. I gràcies a la interacció electromagnètica, juntament amb les lleis de la mecànica quàntica, tenim substàncies, objectes i una matèria raonablement organitzada.

També és la responsable de gairebé tot el que et passa durant el dia. El fregament, la tensió d'una corda, el cop d'una pilota, que la cadira t'aguanti: tot això és electromagnetisme disfressat.

De fet, en sentit estricte, no has tocat mai res. Quan recolzes la mà sobre la taula, els electrons de la teva pell i els de la fusta es repel·leixen abans d'arribar a ajuntar-se. El que anomenem «tocar» és, bàsicament, una repulsió elèctrica a molt curta distància.

\section{Dins del nucli: la forta i la feble}

Ens queden dues forces més. Per trobar-les haurem de mirar dins dels nuclis atòmics.

Però espera un moment.

No acabem de dir que les càrregues del mateix signe es repel·leixen? Aleshores, per què els protons, tots positius, continuen allà atapeïts dins del nucli sense sortir disparats cadascun en una direcció?

D'això se n'encarrega la \textbf{força nuclear forta}. A distàncies diminutes (de l'ordre de \(10^{-15}\) m, la mida d'un nucli) és capaç de vèncer la repulsió elèctrica entre els protons i mantenir units protons i neutrons.

En la seva forma més fonamental, aquesta interacció és la que manté units els quarks que formen aquestes partícules. El que acaba subjectant el nucli és, en realitat, el que sobra d'aquella unió: una mena de resta, com l'olor de la cuina del restaurant de sota. I tot i així és la més intensa de les quatre.

A canvi té un abast ridícul. Separa una mica els protons i deixa de notar-se; per això els nuclis molt grans acaben sent inestables, i per això existeix la radioactivitat.

I encara ens en queda una última: la \textbf{força nuclear feble}.

A diferència de la forta, aquesta no es dedica a mantenir unit res. Intervé en certes transformacions de partícules: per exemple, permet que un neutró es transformi en un protó.

Sembla la menys espectacular de les quatre i té, de bon tros, el pitjor nom. Però sense ella el Sol no es podria encendre, així que més val tractar-la amb respecte.

\section{Uns tipus amb sort}

Veient la quantitat de coses que semblen haver sortit bé, un comença a sospitar que hem tingut una sort de campionat amb això de les forces fonamentals.

Fem recompte.

La gravetat és la força que va fer que la pols i el gas de l'univers primitiu col·lapsessin per formar estrelles, planetes i galàxies. Sense ella, la matèria no s'hauria agrupat i la vida tal com la coneixem no tindria ni on posar-se. A més ens manté enganxats a terra, subjecta l'atmosfera i manté la Terra girant establement al voltant del Sol.

\textbf{Primera carambola.} Que la gravetat tingui exactament la intensitat justa perquè les estrelles cuinin elements durant milers de milions d'anys, els planetes es quedin quietets al seu lloc i tu puguis aixecar-te del sofà sense sortir surant cap al sostre.

Els electrons romanen lligats als àtoms. Altrament, en lloc de taules, arbres i persones, tindríem una sopa de partícules bastant menys acollidora. La força electromagnètica és, a més, moltíssim més intensa que la gravetat: entre dos protons, per exemple, la repulsió elèctrica és aproximadament \(10^{36}\) vegades més gran que la seva atracció gravitatòria.

Aquest número no hi ha manera d'imaginar-lo, així que no ho intentis: és un u seguit de trenta-sis zeros. Si la gravetat jugués a segona regional, l'electromagnetisme estaria en una altra lliga, en un altre país i en un altre esport.

\textbf{Segona carambola.} Que existeixin càrregues positives i negatives que, a la matèria ordinària, es compensin gairebé exactament. Si no fos així, la repulsió elèctrica dominaria bona part del nostre món, i construir una casa amb materials que es repel·leixen entre si seria un problema.

\textit{(Aquí hi ha una mica de trampa, tot sigui dit: precisament perquè l'electromagnetisme és tan brutalment intens, qualsevol excés de càrrega es neutralitza de seguida. No és tant que hàgim tingut sort amb l'equilibri; és que el desequilibri dura poquíssim.)}

Com ja hem comentat, existeix una força, la interacció forta, capaç de mantenir units els nuclis atòmics. Sense ella no tindríem la varietat d'elements químics que coneixem. I n'existeix una altra, la interacció feble, que participa en els processos nuclears que permeten brillar al Sol.

\textbf{Tercera carambola.} Les intensitats de les quatre forces fonamentals, i la resta de constants físiques, tenen valors que permeten l'existència d'àtoms estables, estrelles duradores i una química complexa.

\section{L'ajust fi, o on fa mal}

De fet, acabem de ficar-nos en una de les qüestions més incòmodes (i més fascinants) de la física teòrica actual: el \textbf{problema de l'ajust fi} (\textit{fine-tuning problem}).

La relació entre la força de la gravetat i la nuclear forta, la massa de l'electró davant la del protó, la diferència de massa entre el neutró i el protó... totes aquestes quantitats semblen estar «calibrades» d'una manera molt precisa. Un pèl a l'esquerra i no hi ha química. Un pèl a la dreta i no hi ha estrelles estables.

Arribats aquí se solen sentir tres respostes, i convé saber quina és quina.

La primera: que existeixi alguna teoria més profunda, encara per descobrir, que obligui les constants a valer el que valen. És la que més il·lusió ens fa i la que menys proves té.

La segona: que existeixin moltíssims universos amb constants diferents i que nosaltres estiguem, inevitablement, en un dels pocs on pot haver-hi algú preguntant-s'ho. Té l'inconvenient que serveix per explicar absolutament qualsevol cosa, i una explicació que ho explica tot sol explicar més aviat poc.

La tercera, i la menys popular a les xerrades de divulgació: que la pregunta estigui mal plantejada. Per dir que una cosa és \textit{improbable} cal saber quins altres valors hauria pogut prendre i amb quina probabilitat. I resulta que tenim exactament un univers mesurat i cap altre amb què comparar.

No cal que triïs bàndol avui.

\section{Fabricant forces}

Bé, un moment.

Però no estàvem parlant d'una pilota que es movia? Què hi té a veure tot això amb que algú doni una puntada de peu a una pilota?

Amb la gravetat encara s'entén: una poma cau perquè la Terra l'atrau. Però i les altres tres forces? També tenen alguna cosa a veure amb que es mogui una pilota?

Doncs sí. Encara que sembli estrany, les forces que trobem a la nostra vida quotidiana es poden explicar a partir de les interaccions fonamentals.

I això ens porta a distingir dues parts de la física que farem servir constantment.

La \textbf{cinemàtica} estudia com es mouen els objectes: les seves posicions, velocitats i acceleracions. No es preocupa de per què es mouen; simplement descriu el que fan. És la física fent de notari.

La \textbf{dinàmica}, en canvi, intenta explicar per què canvia aquest moviment. I aquí és on les forces fan la seva aparició.

Però hi ha un petit problema pràctic.

Encara que la força que exerceix el nostre peu sobre la pilota té l'origen, fonamentalment, en interaccions electromagnètiques entre els àtoms, no sembla gaire raonable posar-se a calcular què està fent cada electró per esbrinar fins on arribarà la pilota.

Podríem intentar-ho. Però segurament s'acabaria el partit abans que haguéssim resolt el servei inicial.

Per això, en dinàmica treballem amb forces com el pes, la tensió d'una corda, el fregament o l'empenta que fem amb la mà a una caixa.

No són noves forces fonamentals. Són maneres de descriure, a la nostra escala, els efectes conjunts de moltíssimes interaccions microscòpiques.

I són forces amb totes les de la llei: no necessitem conèixer el comportament de cada partícula per calcular els seus efectes sobre el moviment d'un objecte.

És exactament el mateix que fem quan diem que l'aigua «bull a 100 °C» en lloc de seguir la pista a cada molècula. Es mesura, es prediu i funciona. Que és, al cap i a la fi, l'únic que demanem a un model.

Al capdavall, si dones una puntada de peu a una pilota, la pilota surt disparada.

I això és precisament el que volem explicar.
